% answers_short.tex —— 2002 年·苏豫粤桂卷 · 简短答案速查
% 由 tools/make_short_answers.py 自 answers.tex 的 \gksolhead/\gkanswer 提取，
% 并与 2002/2002_数学真题_苏豫粤桂_答案要点.md 逐题比对后生成，勿手改（改源 MD 后重跑）。

\begin{gknotice}
\textbf{使用说明}：本节只给最终答案便于快速核对，完整推导见解析版。答案为编写者独立推导，\textbf{非官方标准答案}。
\end{gknotice}

\gkpart{选择题}
\begin{center}\begin{tabular}{ccccccccccccc}
\toprule
题号 & 1 & 2 & 3 & 4 & 5 & 6 & 7 & 8 & 9 & 10 & 11 & 12 \\
\midrule
答案 & C & A & D & C & B & C & D & D & C & B & B & C \\
\bottomrule
\end{tabular}\end{center}

\gkpart{填空题}
\begin{itemize}[leftmargin=2.6em]
  \item[\textbf{13.}] $1$。
  \item[\textbf{14.}] $1008$。
  \item[\textbf{15.}] $-\dfrac{\sqrt3}{3}$。
  \item[\textbf{16.}] $\dfrac72$。
\end{itemize}

\gkpart{解答题}
\begin{itemize}[leftmargin=2.6em]
  \item[\textbf{17.}] $(a,b)=(-2,-1)$ 或 $(-4,2)$。
  \item[\textbf{18.}] U/PENDING。
  \item[\textbf{19.}] (1) $V=\dfrac{\sqrt3}{3}a^3$；(2) 按棱锥两侧面沿公共棱围成的内二面角理解，恒大于 $90^\circ$
  \item[\textbf{20.}] (1) $y=x+1$；(2) 四点共圆，圆心 $(-3,6)$，半径 $2\sqrt{10}$
  \item[\textbf{21.}] (1) 连接正三角形三边中点分成四块；一块作正四面体底面，另三块作侧面。另一张纸也按中点分割，选两块作棱柱底面，其余两块按内心至顶点分割并配对、剪角平移成三个矩形侧面；(2) 本构造下正三棱柱体积较大，$V_{\rm prism}=s^3/8>V_{\rm tetra}=\sqrt2s^3/12$；(3) 任意三角形仍按三边中点分成四块，选两块作底面；其余两块由内心至顶点分割，配对并剪角平移为等高矩形侧面，构成直三棱柱
  \item[\textbf{22.}] (1) $a\leqslant2\sqrt b$；(2) $b-1\leqslant a\leqslant2\sqrt b$；(3) $0<a\leqslant b+1$
\end{itemize}
